\section*{Abstract}
\addcontentsline{toc}{section}{Abstract}
Underwater robots need to know which parts of a camera image show divers,
robots, fish, reefs, plants, wrecks, sea floor or open water. We study this as
eight-class semantic segmentation on the SUIM dataset, predicting one class for
every pixel. After auditing the image and mask pairs, we split the official
training pool by duplicate group into \suimTrainImages{} training and
\suimValImages{} validation images and keep the \testImages{} official test
images aside. We compare a U-Net trained from scratch with a DeepLabV3-ResNet50
fine-tuned from pretrained weights. Both models train for \unetEpochs{} epochs
with seed \unetSeed{}, we keep the checkpoint with the highest validation mIoU,
and we score each selected model once on the test set. On the test set,
DeepLabV3 reaches \deeplabCeTestMiou\% mIoU and \deeplabCeTestDice\% mean Dice,
compared with \unetTestMiou\% and \unetTestDice\% for U-Net, at the cost of
\parameterRatio{} times the parameters and \flopRatio{} times the operations
per image. In a controlled experiment, adding a Dice term to the cross-entropy
loss raises plant IoU but lowers robot and wreck IoU on both validation and
test, and its change in validation mIoU is smaller than the variation we
measured between two identical training runs, so the result does not support
our hypothesis. The main remaining errors are plants predicted as reef, which
take \testPlantsAsReef\% of plant pixels on the test set, and mistakes along
object boundaries. Each configuration uses a single seed, which limits the
strength of these comparisons.
